\subsection{R 的开区间之间的同胚}

命题 1. \(\mathbf{R}\) 的所有非空开区间都同胚于 \(\mathbf{R}\) 。先考察有界开区间 \(\mathbf{I} = ]a\) , \(b[(a < b)\) 。对每个 \(x \in \mathbf{I}\) ，令 \(f(x) = -\left(\frac{1}{x - a} + \frac{1}{x - b}\right)\) 。这个函数在 \(\mathbf{I}\) 上连续且严格递增，因为我们已看到 \(\frac{1}{x - b}\) 在 \(]\leftarrow, b[\) 中严格递减，而 \(\frac{1}{x - a}\) 在 \(]a, \rightarrow[\) 中严格递减。由此推出 \(f\) 是 \(\mathbf{I}\) 到一个区间 \(f(\mathbf{I})\) （ \(\mathbf{R}\) 中的）上的同胚（ \(\S 2\) 第 6 目，定理 5）。\(f(\mathbf{I})\) 既不上有界也不下有界；因为若例如 \(f(x) \leqslant c\) 对所有 \(x \in \mathbf{I}\) 成立，则将得出 \(1 - \frac{b - x}{x - a} \leqslant c(b - x)\) （因为 \(b - x > 0\) ）；而当 \(x\) 充分接近 \(b\) 时这就导致矛盾（利用不等式两边这两个有理函数在点 \(b\) 处的连续性）。因而 \(f(\mathbf{I}) = \mathbf{R}\) ，从而每个有界开区间都同胚于 \(\mathbf{R}\) 。设 \(g\) 为 \(f\) 的逆：它把 \(\mathbf{R}\) 的每个无界开区间映到一个区间 \(J\) （含于 \(I\) 且在 \(I\) 中开的）上。由于 \(I\) 在 \(\mathbf{R}\) 中开，\(J\) 在 \(\mathbf{R}\) 中也开；由于它有界，它同胚于 \(\mathbf{R}\) ；这样我们就证明了每个无界开区间都同胚于 \(\mathbf{R}\) 。

注. 为了证明所有有界开区间彼此同胚，只需注意：若 \(a \neq b\) 且 \(a' \neq b'\) ，则存在 R 到其自身的一个形如 \(x \to \alpha x + \beta\) 的同胚（且只有一个），它把 a 映到 \(a'\) ，把 b 映到 \(b'\) ，从而把以 a, b 为端点的开（相应地半开、闭）区间映到以 \(a', b'\) 为端点的开（相应地半开、闭）区间；读者通过计算 \(\alpha\) 与 \(\beta\) 容易验证这一点。

